\chapterwithoutnumber{Заключение}

Такии образом, в рамках данной выпускной работы магистра была реализована система мультиобъектного отслеживания на основе алгоритма DeepSORT с применением нелинейного вильтра предсказания. В ходе работы:

\begin{itemize}

    \item был проведен анализ современных методов детекции и отслеживания для объектов, применяемых в сфере роботизированных ТС;
    \item было приведено обоснования выбора детектора YOLO11n и дальнейшая его тренрировка на наборе данных KITTI;
    \item разработан алгоритм МОТ на основе DeepSORT с использованием нелинейного фильтра предсказания положения объекта;
    \item разраобтан алгоритм определения НВО;
    \item проведено тестирование камера-детектора YOLO11n и разработанного алгоритма МОТ;
    \item проведен анализ результатов тестирования, а также сравнение разработанного алгоритма с аналогами и доказана его конкурентноспособность.

\end{itemize}

На основании вышеизложенного можно сделать следующие выводы:

\begin{itemize}

    \item разработанная система превосходит своих конкурентов в сценах с плотным трафиком и наличием большого количества элементов, имеющих нелинейные траектории;
    \item разработанная система не показывает ухудшения в сценах с небольшим трафиком в хороших погодных условиях и с линейным поведением объектов;
    \item точность работы системы зависит от погодных условий и освещения, но не опускается ниже требуемой нормы, что остается приемлимым для безопасного движения. 

\end{itemize}

Исходя из сказанного, можно сделать вывод о работоспособности и возможности конкурирования на рынке разработанной системы мультиобъектного трекинга с примнением нелинейного фильтра предсказания.